\section{Pendefinisian Masalah Data}
Tahap awal dalam pengembangan model \textit{machine learning} adalah mendefinisikan karakteristik data dan batasan masalah untuk memastikan model yang dibangun relevan dengan tujuan penelitian, yaitu mendeteksi perilaku penyewa yang berpotensi melakukan pelanggaran di Rusunami The Jarrdin Cihampelas (RTJC).

\subsection{Tujuan Penggunaan Data}
Data yang dikumpulkan bertujuan untuk membangun model klasifikasi biner (\textit{binary classification}) yang dapat memisahkan penyewa ke dalam dua kelas:
\begin{itemize}
    \item \textbf{Kelas normal (0):} penyewa yang mematuhi aturan administrasi dan tidak memiliki catatan pelanggaran.
    \item \textbf{Kelas tidak normal (1):} penyewa yang terindikasi melakukan pelanggaran, yang ditandai dengan tidak kosongnya kolom komentar pelanggaran pada basis data.
\end{itemize}

\subsection{Kriteria dan Pelabelan Data}
Berdasarkan tinjauan pustaka pada Bab~\ref{chap:teori} mengenai \textit{Supervised Learning}, proses pelabelan data (\textit{labeling}) dilakukan secara objektif berdasarkan fakta historis yang tersimpan. Sesuai dengan kode pengembangan model, pelabelan dilakukan dengan logika sebagai berikut:

\[
Label(x) =
\begin{cases}
0, & \text{jika } komentar \in \{\text{``'', ``nan'', ``-'', ``0'', ``tidak ada''}\} \\
1, & \text{lainnya (terdapat catatan pelanggaran)}
\end{cases}
\]
Pendekatan ini dipilih untuk mengubah data teks komentar yang tidak terstruktur menjadi target biner yang dapat dipelajari oleh model \textit{machine learning}.

\subsection{Tantangan Ketidakseimbangan Kelas (\textit{Class Imbalance})}
Berdasarkan eksplorasi awal pada dataset yang terdiri dari 9.989 data transaksi, ditemukan bahwa proporsi data pelanggaran (label 1) hanya berjumlah 746 data, atau sekitar 7,47\% dari total populasi.

Ketidakseimbangan ini menjadi tantangan utama karena algoritma standar cenderung bias terhadap kelas mayoritas. Oleh karena itu, pengembangan model memprioritaskan algoritma yang mendukung penanganan \textit{imbalance} seperti \textit{Balanced Random Forest} dan teknik \textit{cost-sensitive learning} (parameter \texttt{class\_weight='balanced'}).